% ============================================================
% DOCUMENTCLASS
% ============================================================
\documentclass[12pt,oneside]{book}

% ============================================================
% METADATOS
% ============================================================
\newcommand{\universidad}{Universidad de Buenos Aires (UBA)}
\newcommand{\materia}{Derecho Civil}
\newcommand{\titulo}{Apuntes de Derecho Civil: Constitucionalización del Derecho Privado y Título Preliminar del CCyC}
\newcommand{\autor}{Isaac Edgar Camacho Ocampo}
\newcommand{\anio}{2025}

% ============================================================
% PAQUETES
% ============================================================
\usepackage[utf8]{inputenc}
\usepackage[T1]{fontenc}
\usepackage[spanish,es-nodecimaldot]{babel}

\usepackage[
    top=2.5cm,
    bottom=2.5cm,
    left=3cm,
    right=2.5cm,
    headheight=15pt
]{geometry}

\usepackage{amsmath}
\usepackage{amssymb}
\usepackage{mathtools}

\usepackage{graphicx}
\usepackage{float}
\usepackage{caption}
\usepackage{subcaption}

\usepackage{booktabs}
\usepackage{array}
\usepackage{longtable}
\usepackage{tabularx}

\usepackage{enumitem}

\usepackage{xcolor}
\usepackage[most]{tcolorbox}

\usepackage{fancyhdr}

\usepackage{ragged2e}

\usepackage[
    colorlinks=true,
    linkcolor=blue!50!black,
    citecolor=green!40!black,
    urlcolor=blue!60!black,
    pdftitle={\titulo},
    pdfauthor={\autor},
    pdfsubject={Apuntes universitarios},
    bookmarks=true
]{hyperref}

% ============================================================
% CONFIGURACIÓN
% ============================================================
\graphicspath{
    {./img/}
    {./graficos/}
    {./figuras/}
}

\setcounter{tocdepth}{2}

\setlength{\parindent}{0pt}
\setlength{\parskip}{0.5em}

\setlist[itemize]{
    topsep=0.3em,
    itemsep=0.2em
}

\setlist[enumerate]{
    topsep=0.3em,
    itemsep=0.2em
}

\clubpenalty=10000
\widowpenalty=10000

% ============================================================
% COMANDOS
% ============================================================
\newcommand{\abs}[1]{\left|#1\right|}
\newcommand{\norm}[1]{\left\lVert#1\right\rVert}
\newcommand{\vect}[1]{\mathbf{#1}}
\newcommand{\deriv}[2]{\frac{d#1}{d#2}}
\newcommand{\pderiv}[2]{\frac{\partial#1}{\partial#2}}

% ============================================================
% ENTORNOS / CAJAS
% ============================================================
\newtcolorbox{definition}[1]{
    enhanced,
    breakable,
    title={Definición: #1},
    fonttitle=\bfseries,
    colback=blue!3,
    colframe=blue!50!black,
    boxrule=0.8pt,
    arc=2mm
}

\newtcolorbox{important}{
    enhanced,
    breakable,
    title={Importante},
    fonttitle=\bfseries,
    colback=yellow!8,
    colframe=orange!70!black,
    boxrule=0.8pt,
    arc=2mm
}

\newtcolorbox{example}[1]{
    enhanced,
    breakable,
    title={Ejemplo: #1},
    fonttitle=\bfseries,
    colback=green!4,
    colframe=green!40!black,
    boxrule=0.8pt,
    arc=2mm
}

\newtcolorbox{formula}[1]{
    enhanced,
    breakable,
    title={Fórmula / Regla: #1},
    fonttitle=\bfseries,
    colback=purple!4,
    colframe=purple!50!black,
    boxrule=0.8pt,
    arc=2mm
}

\newtcolorbox{exercise}[1]{
    enhanced,
    breakable,
    title={Ejercicio: #1},
    fonttitle=\bfseries,
    colback=gray!5,
    colframe=gray!60!black,
    boxrule=0.8pt,
    arc=2mm
}

\newtcolorbox{note}{
    enhanced,
    breakable,
    title={Nota},
    fonttitle=\bfseries,
    colback=white,
    colframe=gray!50,
    boxrule=0.6pt,
    arc=2mm
}

% ============================================================
% CONFIGURACIÓN FINAL (encabezados y pies)
% ============================================================
\pagestyle{fancy}
\fancyhf{}

\fancyhead[L]{\nouppercase{\leftmark}}
\fancyhead[R]{\materia}

\fancyfoot[C]{\thepage}

\renewcommand{\headrulewidth}{0.4pt}
\renewcommand{\footrulewidth}{0pt}

\fancypagestyle{plain}{
    \fancyhf{}
    \fancyfoot[C]{\thepage}
    \renewcommand{\headrulewidth}{0pt}
}

% ============================================================
% DOCUMENT
% ============================================================
\begin{document}

% ------------------------------------------------------------
% PORTADA
% ------------------------------------------------------------
\begin{titlepage}

\thispagestyle{empty}

\begin{center}

\vspace*{1cm}

{\large \textsc{\universidad}}

\vspace{3cm}

{\Huge \textbf{\materia}}

\vspace{1cm}

{\LARGE \titulo}

\vspace{0.8cm}

\textit{Comprender los conceptos es el primer paso para convertir el estudio en conocimiento.}

\vspace{1.2cm}

\rule{0.70\textwidth}{0.5pt}

\vspace{0.5cm}

{\large Apuntes de estudio}

\vspace{0.5cm}

\rule{0.70\textwidth}{0.5pt}

\vfill

{\large \autor}

\vspace{0.3cm}

{\large \anio}

\end{center}

\vfill

\begin{flushleft}
\textit{Este es un modesto aporte a los estudiantes de ``Derecho Civil'' no pretende sustituir el material oficial de la materia.}
\end{flushleft}

\end{titlepage}

% ------------------------------------------------------------
% ÍNDICE
% ------------------------------------------------------------
\tableofcontents

% ------------------------------------------------------------
% CONTENIDO
% ------------------------------------------------------------
\chapter{Constitucionalización del derecho privado y el Título Preliminar del Código Civil y Comercial}

En esta clase se profundiza el tema de la constitucionalización del derecho privado y el nuevo Código Civil y Comercial (CCyC), con particular referencia a su Título Preliminar. Los temas tratados son:

\begin{itemize}
    \item Constitucionalización del derecho privado.
    \item Las tres vías de constitucionalización (reforma, hermenéutica y aplicación directa).
    \item Ventajas y riesgos de la constitucionalización.
    \item El Título Preliminar del Código Civil y Comercial argentino (2014/2015).
    \item Unificación de la materia civil y comercial.
    \item Principios fundamentales del derecho civil: justicia, buena fe y equidad.
\end{itemize}

\begin{note}
\textbf{Relevancia profesional.} Comprender la constitucionalización del derecho privado resulta esencial para cualquier abogado, escribano u operador jurídico: permite argumentar con los estándares más altos del ordenamiento (Constitución y tratados de derechos humanos), anticipar cómo un juez interpretará un contrato, una herencia o una relación de familia, e identificar cuándo una norma puede ser cuestionada por violar principios superiores. Los principios de buena fe, equidad y dignidad de la persona son herramientas cotidianas en el litigio, la asesoría y la redacción de contratos. Además, tienen aplicación directa en la vida cotidiana: cada vez que se firma un contrato o se tiene un problema con un prestador de servicios, la buena fe, la equidad y la dignidad de la persona son los principios que un juez aplicará si una cláusula resulta abusiva o si los derechos de una persona son vulnerados.
\end{note}

\begin{example}{Analogía explicativa de la lógica de la clase}
Imaginemos que la Constitución es la estructura de un edificio: define dónde pueden ir las paredes y cuánto peso pueden soportar. El Código Civil y Comercial es el interior del edificio (habitaciones, puertas, contratos, familia, herencias), pero su diseño debe respetar esa estructura madre. El Título Preliminar es el plano que explica cómo leer todo el edificio: qué fuentes son válidas, cómo interpretar las normas y qué principios prevalecen. Sin ese plano, cada abogado leería el edificio a su manera.
\end{example}

\section{Constitucionalización del derecho privado}

\subsection{¿Qué significa constitucionalizar el derecho privado?}

Cuando se piensa el tema de la constitucionalización, se hace referencia a qué influencia tiene la Constitución Nacional —y todos los tratados de derechos humanos— sobre el derecho privado. Este tema cobra relevancia, sobre todo, luego de la reforma constitucional de 1994, que incorporó con jerarquía constitucional los tratados internacionales de derechos humanos.

\begin{important}
La reforma constitucional de 1994 otorgó jerarquía constitucional a los tratados internacionales de derechos humanos (art. 75 inc. 22 CN). Esta jerarquización proyecta efectos de distinto tipo sobre todas las ramas jurídicas, incluyendo la del derecho civil.
\end{important}

Tal jerarquización proyecta efectos de distinto tipo sobre todas las ramas jurídicas, incluyendo la rama del derecho civil.

\subsection{Las tres vías de constitucionalización (según Corral Talciani)}

El jurista chileno Corral Talciani explica que la constitucionalización opera por \textbf{tres vías}.

\subsubsection{La vía reformadora}

A partir de los tratados de derechos humanos y sus disposiciones, se generan cambios en el ámbito legislativo.

\begin{example}{Vía reformadora — Convención sobre los Derechos de las Personas con Discapacidad}
En la Argentina, cuando en el año 2008 se ratificó la Convención sobre los Derechos de las Personas con Discapacidad —en particular su artículo 12, que señala una serie de exigencias en torno a cómo dar regulación jurídica a la temática vinculada con la capacidad de ejercicio y la capacidad jurídica en general—, Argentina receptó este tratado y esto llevó a que se planteara la necesidad de revisar su propia legislación. En el año 2010 se sancionó una ley que modificó parcialmente —en aquel momento— el Código Civil de Vélez Sársfield en lo referido a cómo debían ser las sentencias en materia de capacidad de ejercicio.
\end{example}

Si se piensa en la Convención sobre los Derechos del Niño, el interés superior del niño como principio y el derecho del niño a ser oído también fueron generando cambios legislativos sucesivos en el país hasta el Código Civil y Comercial. El Código Civil y Comercial es muy sensible a este tema: recepta las indicaciones de la Convención sobre los Derechos del Niño y de distintas convenciones en materia de derechos humanos en sus disposiciones.

\subsubsection{La vía hermenéutica}

La vía hermenéutica refiere a la forma de interpretar. Los principios y los tratados de derechos humanos están allí para interpretar el alcance de la ley, como criterio al momento de juzgar.

\begin{formula}{Artículo 2° del CCyC --- Interpretación}
La ley debe ser interpretada teniendo en cuenta sus palabras, sus finalidades, las leyes análogas, las disposiciones que surgen de los tratados sobre derechos humanos, los principios y los valores jurídicos, de modo coherente con todo el ordenamiento.
\end{formula}

Este artículo recoge en el texto del código la importancia de los tratados de derechos humanos como guía interpretativa, reflejando la vía hermenéutica de constitucionalización.

\subsubsection{La vía de aplicación directa}

En algunos casos se asiste a una aplicación directa de los tratados de derechos humanos y los principios constitucionales. La Constitución del año 1994 incorporó normas sobre el consumidor y sobre habeas data.

Estas normas se aplican a través de legislación específica —en el caso del consumidor, la Ley de Defensa del Consumidor; en el caso de los datos personales, la Ley de Protección de Datos Personales—, pero, llegado a ciertos casos, puede haber una aplicación directa de las normas constitucionales. Esto es más raro porque generalmente el texto constitucional aparece mediado, concretado a través de la legislación y la tarea propia del legislativo, pero puede darse el caso en que directamente se apele a un principio constitucional, como por ejemplo la dignidad.

\begin{important}
\textbf{Caso emblemático (2019):} la Corte Suprema argentina recurrió al principio de la dignidad y del derecho a la vida para modificar —en favor de una persona con discapacidad— el orden en que se cobran los créditos en el régimen de privilegios de la quiebra (Ley de Concursos y Quiebras). Se trató de un caso de aplicación directa de un principio constitucional para modificar una ley, situación reconocida como excepcional.
\end{important}

\section{Ventajas y riesgos de la constitucionalización}

\subsection{Ventajas}
% TODO: contenido incompleto o ambiguo en las notas originales — el apunte atribuye la enumeración de ventajas «al profesor Daniela Herrera», con discordancia de género entre el artículo y el nombre propio; se conserva la referencia tal como figura en el material de origen.

Siguiendo a la profesora Daniela Herrera, la constitucionalización ha traído importantes ventajas para el derecho privado:

\begin{longtable}{
    >{\raggedright\arraybackslash}p{0.24\textwidth}
    >{\raggedright\arraybackslash}p{0.38\textwidth}
    >{\raggedright\arraybackslash}p{0.28\textwidth}
}
\toprule
\textbf{Ventaja} & \textbf{Descripción} & \textbf{Reflejo en el CCyC} \\
\midrule
\endhead
Centralidad de la persona y su dignidad &
Se coloca a la persona humana en el centro del ordenamiento jurídico, resaltando su dignidad. &
Art. 51 (dignidad de la persona humana), Art. 19 (inicio de la existencia), Art. 57 (protección del embrión) \\
Re-materialización del derecho &
Incorporación de contenidos sustanciales frente a la idea de un derecho meramente procedimental o formal. &
Todo el Título Preliminar y los artículos de la parte general sobre persona y familia \\
Rehabilitación de la dimensión práctica o valorativa &
Se recupera la importancia de la argumentación y de los principios en la tarea de juzgar. &
Art. 2 (interpretación), Art. 3 (deber de decidir de modo razonablemente fundado) \\
La Constitución como fuente del derecho &
La Constitución y sus principios y valores se transforman en una fuente del derecho aplicable. &
Art. 1 (fuentes del derecho: Constitución y tratados junto a la ley) \\
\bottomrule
\end{longtable}

\begin{formula}{Artículo 1° del CCyC --- Fuentes y aplicación}
Los casos que este Código rige deben ser resueltos según las leyes que resulten aplicables, conforme con la Constitución Nacional y los tratados de derechos humanos en los que la República sea parte. A tal efecto, se tendrá en cuenta la finalidad de la norma. Los usos, prácticas y costumbres son vinculantes cuando las leyes o los interesados se refieren a ellos o en situaciones no regladas legalmente, siempre que no sean contrarios a derecho.
\end{formula}

\subsection{Riesgos de la constitucionalización}

Desde el punto de vista de la filosofía jurídica, quienes sostienen una posición iusnaturalista entienden que se corre el riesgo de pasar de un positivismo legalista a un positivismo constitucionalista, perdiendo la dimensión de principios superiores al orden jurídico que iluminan e inciden en la valoración del ordenamiento.

\begin{longtable}{
    >{\raggedright\arraybackslash}p{0.30\textwidth}
    >{\raggedright\arraybackslash}p{0.62\textwidth}
}
\toprule
\textbf{Riesgo} & \textbf{Descripción} \\
\midrule
\endhead
Pérdida de especificidad del derecho civil &
La expansión del derecho constitucional puede afectar las metodologías propias de las otras ramas. El derecho civil tiene principios jurídicos sectoriales propios —como las reglas para calcular una indemnización por daños y perjuicios— que pueden perderse cuando todo se constitucionaliza. \\
Positivismo judicial &
El juez que, al aplicar la Constitución, se aparta excesivamente de la finalidad de la ley estaría determinando qué es lo obligatorio, alterando la división de poderes y erosionando la construcción republicana. \\
Inseguridad jurídica &
Si en las relaciones civiles no se supiera a qué atenerse —cómo serán los derechos y obligaciones de cada uno—, la enorme inseguridad jurídica resultante paralizaría la vida económica y social, generaría mayores costos y crearía problemas nuevos. \\
Desborde de fuentes del derecho &
Por la pérdida de relevancia de la ley, aparecen como fuentes los fallos de tribunales internacionales, recomendaciones y convenciones de otros sistemas jurídicos, lo que genera una dificultad para determinar cuál es el derecho aplicable ante una situación determinada. \\
Relativismo filosófico &
A veces se piensa que las posiciones de constitucionalización del derecho privado sólo son admisibles en el marco de una construcción filosófica relativista. Sin embargo, desde el iusnaturalismo se sostiene que es posible conocer —a partir de la racionalidad humana— principios fundamentales: que la vida es un bien y que quitarle la vida a una persona es un mal. \\
\bottomrule
\end{longtable}

Es ciertamente un profundo debate sobre el juez, la jurisprudencia como fuente, en qué medida el juez crea el derecho o lo está aplicando, y cuál es su tarea ante un caso concreto cuando existe ley que regula esa situación.

\section{El Título Preliminar del Código Civil y Comercial}

\subsection{Importancia del Título Preliminar en el proceso de constitucionalización}

Al mirar el Código Civil y Comercial se descubre la importancia del Título Preliminar en este proceso de constitucionalización, porque justamente en el artículo 1 se señala la importancia de la Constitución y los tratados internacionales de derechos humanos como pautas para aplicar a los casos regidos por el código.

\begin{important}
El Título Preliminar tiene la función de ofrecer un núcleo de significaciones general para todo el código. Es como si se quisiera señalar que en torno al Título Preliminar se encuentran principios fundamentales que influyen sobre todo el resto de los libros del código.
\end{important}

\subsection{Estructura del Título Preliminar}

El Título Preliminar del Código Civil y Comercial se divide en cuatro capítulos:

\begin{longtable}{
    >{\raggedright\arraybackslash}p{0.08\textwidth}
    >{\raggedright\arraybackslash}p{0.22\textwidth}
    >{\raggedright\arraybackslash}p{0.40\textwidth}
    >{\raggedright\arraybackslash}p{0.18\textwidth}
}
\toprule
\textbf{Cap.} & \textbf{Título} & \textbf{Contenido principal} & \textbf{Artículos} \\
\midrule
\endhead
1 & El derecho & Fuentes y aplicación. Distinción entre derecho y ley. & Arts. 1, 2 y 3 \\
2 & La ley & Ámbito de vigencia temporal y espacial. Efectos de la ley en relación al tiempo. & Arts. 4 a 8 \\
3 & Ejercicio de los derechos & Principio del abuso del derecho y principio de buena fe. & Arts. 9 a 14 \\
4 & Derechos y bienes & Clasificación de bienes jurídicos, cosas, patrimonio, cuerpo humano. & Arts. 15 a 18 \\
\bottomrule
\end{longtable}

El Título Preliminar consta de 18 artículos. En la elevación del proyecto del código se señala que estos artículos proyectan sus efectos sobre los seis libros que componen el resto del Código Civil y Comercial.

\subsection{La distinción entre derecho y ley}

Es significativo que se haya distinguido entre «derecho» y «ley» en los fundamentos de elaboración del proyecto. En los fundamentos se llama la atención sobre esta distinción, señalando que el derecho es más que la ley.

Cuando se lee el contenido del capítulo primero —que comprende los tres primeros artículos— se advierte que:

\begin{itemize}
    \item El artículo 1 habla de las fuentes (Constitución, tratados, ley, costumbre).
    \item El artículo 2 habla de la interpretación de la ley.
    \item El artículo 3 establece que en todo caso debe emitirse una decisión razonablemente fundada.
\end{itemize}

La importancia de «el derecho», entendido como algo que supera a la ley, está enfatizada en el Título Preliminar, aunque en la práctica la ley sigue siendo la principal de las fuentes, según se aprecia en el capítulo 2 del mismo título.

\section{Unificación de la materia civil y comercial}

\subsection{Antecedentes y alcance de la unificación}

En la Argentina existieron distintas iniciativas de unificación de la materia civil y comercial. Finalmente, el Código Civil y Comercial lo concreta, aunque —según señalan los autores que han estudiado el tema— se trata de una unificación parcial.

Se incorporan al código unificado la teoría general de las condiciones y los contratos, con las instituciones comunes, y ciertos temas en materia de actos jurídicos y de contabilidad y registros contables.

\subsection{Materias que permanecen fuera del CCyC}

Ciertos temas quedan fuera del código y continúan siendo regulados por sus propias leyes especiales:

\begin{itemize}
    \item Las sociedades comerciales (Ley General de Sociedades).
    \item Las quiebras y concursos (Ley de Concursos y Quiebras).
    \item La defensa de la competencia.
    \item La empresa como institución y el empresario no están incluidos ni regulados específicamente en el CCyC.
\end{itemize}

Tampoco hay en el CCyC una definición de cuál es la materia comercial. Hay que recordar que quedó derogado el Código de Comercio, que contenía un artículo específico sobre acto de comercio, la definición de acto de comercio y de comerciante.

\subsection{Contenido comercial incorporado al CCyC}

Lo que sí aparece en el CCyC con contenido originalmente comercial incluye:

\begin{itemize}
    \item La contabilidad y la rendición de cuentas.
    \item Toda la teoría de la representación (arts. 358 y siguientes).
    \item La parte especial de contratos, con muchos contratos comerciales típicos.
    \item Las relaciones de consumo.
    \item Las reglas de interpretación de contratos que estaban antes en el Código de Comercio, ahora incorporadas al CCyC en la parte de contratos.
    \item El valor de los usos y las costumbres: hay una mayor presencia de los usos y las costumbres como fuente en distintos artículos, especialmente en materia de contratos.
    \item Los títulos de crédito.
\end{itemize}

\section{Principios fundamentales del derecho civil}

\subsection{La justicia como principio central}

En definitiva, el derecho civil tiene como gran tema de fondo la justicia —si se quiere, la justicia de «dar a cada uno lo suyo»—. La justicia se presenta como la virtud que ordena la relación con los demás:

\begin{longtable}{
    >{\raggedright\arraybackslash}p{0.22\textwidth}
    >{\raggedright\arraybackslash}p{0.42\textwidth}
    >{\raggedright\arraybackslash}p{0.24\textwidth}
}
\toprule
\textbf{Tipo de justicia} & \textbf{Descripción} & \textbf{Campo de aplicación} \\
\midrule
\endhead
Justicia legal & Del todo hacia las partes: obliga a los miembros de la comunidad hacia el bien común. & Derecho público / constitucional \\
Justicia distributiva & Del todo hacia la parte: la distribución de cargas y beneficios desde la comunidad hacia los individuos. & Derecho administrativo / tributario \\
Justicia conmutativa & Entre las partes: busca la igualdad en los intercambios, la reciprocidad. Es la justicia propia de las relaciones entre particulares. & Derecho civil / contratos / daños \\
\bottomrule
\end{longtable}

El derecho civil se mueve en el campo de la justicia conmutativa, en lo que se ha llamado la «reciprocidad de los cambios». Existen estudios muy a fondo del profesor Daniel Alioto sobre la importancia de la reciprocidad en los cambios como principio fundamental del derecho civil.

En el centro está la persona humana con su dignidad, resaltada en el artículo 51 del CCyC. La persona está protegida desde la concepción hasta la muerte natural —artículo 19—, que señala el inicio de la existencia de la persona conforme a los tratados de derechos humanos.

\begin{formula}{Artículo 51 CCyC --- Inviolabilidad de la persona humana}
La persona humana es inviolable y en cualquier circunstancia tiene derecho al reconocimiento y respeto de su dignidad.
\end{formula}

\begin{formula}{Artículo 19 CCyC --- Comienzo de la existencia}
La existencia de la persona humana comienza con la concepción.
\end{formula}

A su vez, el derecho civil, como toda regulación jurídica, también busca el bien común político y sirve a la realización de un bien que trasciende los bienes individuales: el conjunto de condiciones que permiten que las personas, las familias y las sociedades alcancen su perfección, participando de un bien participado que se llama bien común.

\subsection{El principio de buena fe}

Un principio fundamental es el principio de la buena fe. Los derechos deben ser ejercidos de buena fe.

\begin{formula}{Artículo 9° CCyC --- Principio de buena fe}
Los derechos deben ser ejercidos de buena fe.
\end{formula}

Este principio se incorpora en el Título Preliminar y proyecta sus efectos sobre todo el código. La buena fe tiene dos grandes dimensiones:

\begin{longtable}{
    >{\raggedright\arraybackslash}p{0.20\textwidth}
    >{\raggedright\arraybackslash}p{0.36\textwidth}
    >{\raggedright\arraybackslash}p{0.32\textwidth}
}
\toprule
\textbf{Dimensión} & \textbf{Descripción} & \textbf{Ejemplo} \\
\midrule
\endhead
Buena fe objetiva & La conformidad con ciertas reglas; un obrar que sea conforme a estándares objetivos de comportamiento. & El poseedor es de buena fe cuando no conoce —y no puede conocer— que carece de derecho. Por ejemplo, quien compra algo que tenía un vicio en la cadena de transmisión sin saberlo. \\
Buena fe subjetiva & Un deber de lealtad. Se vincula con cómo se ejecutan y desenvuelven las partes en los contratos: con honestidad, sin ocultar ni retener información relevante. & El deber de honestidad entre contratantes durante las negociaciones y la ejecución del contrato. Una parte no puede retener información que la otra necesita para tomar decisiones. \\
\bottomrule
\end{longtable}

La buena fe subjetiva sirve como un deber de honestidad y lealtad, como una norma de comportamiento de conducta en las relaciones entre particulares.

\subsection{El principio de equidad}

Como gran principio también aparece el principio de la equidad. La equidad se vincula con la justicia en el caso concreto: es la clásica «rectificación de lo justo legal en el caso concreto», según la formulación aristotélica.

La noción de equidad se encuentra ya desde la Ley 17.711 y el Código Civil anterior. En el nuevo Código Civil y Comercial se encuentra sobre todo en los artículos 332 y 1742.

\begin{formula}{Artículo 332 CCyC --- Lesión}
Puede demandarse la nulidad o la modificación de los actos jurídicos cuando una de las partes, explotando la necesidad, debilidad síquica o inexperiencia de la otra, obtuviera por medio de ellos una ventaja patrimonial evidentemente desproporcionada y sin justificación. [\ldots] El juez debe apreciar subjetivamente las circunstancias [\ldots] y los jueces pueden hacer reajustes equitativos cuando se presenta una situación desproporcionada, buscando la corrección de la situación injusta que se presenta en el caso concreto.
\end{formula}

\begin{formula}{Artículo 1742 CCyC --- Atenuación de la responsabilidad}
El juez, al fijar la indemnización, puede atenuarla si es equitativo en función del patrimonio del deudor, la situación personal de la víctima y las circunstancias del hecho. Esta facultad no es aplicable en caso de dolo del responsable.
\end{formula}

En ambos casos, los jueces pueden hacer rectificaciones o reajustes equitativos, buscando la corrección de la situación injusta que se presenta en el caso concreto.

% ------------------------------------------------------------
% CORNELL
% ------------------------------------------------------------
\chapter*{Cuadro Cornell de repaso}
\addcontentsline{toc}{chapter}{Cuadro Cornell de repaso}

Esta sección organiza los conceptos clave de la clase según el método Cornell: preguntas disparadoras a la izquierda, respuestas o explicaciones a la derecha, y una síntesis integradora al pie.

\begin{longtable}{
    >{\raggedright\arraybackslash}p{0.30\textwidth}
    >{\raggedright\arraybackslash}p{0.62\textwidth}
}
\toprule
\textbf{Preguntas / palabras clave} & \textbf{Notas / respuestas} \\
\midrule
\endhead

\textbf{¿Qué es la constitucionalización del derecho privado?} &
Es la influencia que la Constitución Nacional y los tratados de derechos humanos ejercen sobre el derecho privado. Se intensificó con la reforma constitucional de 1994, que otorgó jerarquía constitucional a los tratados internacionales de derechos humanos (art. 75 inc. 22 CN). \\

\textbf{¿Cuáles son las tres vías de constitucionalización según Corral Talciani?} &
1) Vía reformadora: los tratados generan cambios legislativos (ej.: reforma sobre capacidad jurídica en 2010 por la Convención sobre Discapacidad). 2) Vía hermenéutica: los principios y tratados sirven para interpretar la ley (art. 2 CCyC). 3) Vía de aplicación directa: excepcionalmente, se aplica un principio constitucional directamente a un caso (ej.: dignidad en materia de quiebras, 2019). \\

\textbf{¿Qué ventajas trajo la constitucionalización?} &
Centralidad de la persona y su dignidad, re-materialización del derecho (contenidos sustanciales sobre forma), rehabilitación de la argumentación y los principios en la interpretación, y la incorporación de la Constitución como fuente formal del derecho (art. 1 CCyC). \\

\textbf{¿Cuáles son los riesgos de la constitucionalización?} &
Pérdida de especificidad del derecho civil, positivismo judicial (el juez se aparta excesivamente de la ley), inseguridad jurídica, desborde de las fuentes del derecho (fallos internacionales, recomendaciones, etc.) y el riesgo de caer en un relativismo filosófico. \\

\textbf{¿Qué es el Título Preliminar y cuál es su función?} &
Son los 18 artículos iniciales del CCyC, organizados en cuatro capítulos (el derecho, la ley, el ejercicio de los derechos, los derechos y los bienes). Su función es ofrecer un núcleo de significaciones generales que proyecta sus efectos sobre todos los libros del código. \\

\textbf{¿Qué unificó el CCyC y qué quedó fuera?} &
Unificó la teoría general de los contratos, los actos jurídicos, la contabilidad y la representación. Quedaron fuera: las sociedades, las quiebras, la defensa de la competencia. No hay definición de «materia comercial» ni de «acto de comercio» en el código unificado. \\

\textbf{¿Cuál es el campo de la justicia en el derecho civil?} &
El derecho civil opera en el campo de la justicia conmutativa: la que busca la igualdad en los intercambios entre particulares, la reciprocidad. A diferencia de la justicia legal (del individuo hacia el todo) y la distributiva (del todo hacia la parte). \\

\textbf{¿Cuáles son las dos dimensiones de la buena fe?} &
La buena fe objetiva implica actuar conforme a reglas estándares de conducta (ej.: el poseedor de buena fe no conoce ni puede conocer su falta de derecho). La buena fe subjetiva implica un deber de lealtad y honestidad en la ejecución de los contratos, sin ocultar información relevante. \\

\textbf{¿Qué es la equidad y cómo la aplican los jueces?} &
La equidad es la rectificación de lo justo legal en el caso concreto: ajustar la norma general a las particularidades del caso para evitar una solución injusta. Los arts. 332 (lesión) y 1742 (atenuación de la indemnización) habilitan a los jueces a hacer reajustes equitativos cuando la aplicación mecánica de la ley produciría un resultado desproporcionado. \\

\midrule

\multicolumn{2}{p{0.95\textwidth}}{
\textbf{Resumen:} La clase expuso cómo la reforma constitucional argentina de 1994 transformó el modo en que debe leerse, interpretarse y aplicarse el Código Civil y Comercial. La constitucionalización opera por tres vías —reforma legislativa, hermenéutica y aplicación directa—, y aunque trae ventajas indudables (centralidad de la persona, re-materialización del derecho, la Constitución como fuente), también encierra riesgos: pérdida de especificidad del derecho civil, positivismo judicial e inseguridad jurídica. El Título Preliminar del CCyC (18 artículos, aprobado en 2014 y vigente desde 2015) es la clave de bóveda de este proceso: fija las fuentes del derecho, los criterios de interpretación y los principios fundamentales —buena fe, dignidad, equidad— que proyectan sus efectos sobre los seis libros restantes del código. En cuanto a la unificación civil y comercial, el CCyC integró contratos, representación y relaciones de consumo, pero dejó fuera a las sociedades, las quiebras y la empresa como institución. Finalmente, los principios de justicia conmutativa, buena fe (objetiva y subjetiva) y equidad son las columnas vertebrales del derecho civil: determinan cómo se ejercen los derechos, cómo se interpretan los contratos y cómo los jueces corrigen las situaciones injustas en el caso concreto.
} \\

\bottomrule
\end{longtable}

\end{document}
